\chapter{产品与服务}

\section{产品定位与核心理念}

\subsection{极简工具定位：拍照即诊断、诊断即给方案}

「慧农识病」不做功能齐全的农业平台，而是一款\hlgreen{「拍照即诊断、诊断即给方案」的极简工具}。产品只有一个核心交互主线：打开——拍照——上传——得到诊断结果与防治方案。所有功能设计围绕「让 60 岁老农也能在 30 秒内上手」这一目标展开。

\subsection{全闭环设计：识别—防治—预警完整链路}

产品并非孤立的「识别器」，而是打通「\hlgold{识别——防治——预警}」的完整价值闭环：

\begin{itemize}[leftmargin=1.6em]
  \item \hlgreen{识别}：拍照上传，AI 识别病虫害种类并输出置信度；
  \item \hlgreen{防治}：识别结果自动匹配精准防治方案（用药类型、用量、施用时机）；
  \item \hlgreen{预警}：基于地理定位，主动推送本地当季高发病虫害预防提示。
\end{itemize}

\section{核心功能模块}

\subsection{AI 拍照识别诊断}

用户拍摄作物叶片、果实或虫体照片，系统调用端侧轻量化模型完成病虫害种类识别，并以\hlgold{置信度条}直观展示 Top-3 候选结果，降低「识别单一答案」带来的误判风险。

\subsection{精准防治方案匹配}

识别结果自动关联权威防治知识库，输出「用药类型、用药用量、施用时机、注意事项」一体化建议。方案通过多属性决策模型综合药效、成本、毒性、窗口期等维度，给出 Top-3 推荐并按综合得分排序。

\subsection{病虫害知识库搜索}

作为识别功能的补充，农户可主动搜索已知病虫害的防治方法，支持以「作物＋症状描述＋图片」多模态检索，满足「已经知道是什么病，只想查怎么治」的场景。

\subsection{本地化预警推送}

结合用户地理位置，推送当地、当季高发病虫害的预防预警，提醒农户「提前预防」而非「事后补救」，真正抓住防治窗口期。

\section{产品核心特色}

\subsection{极致轻量化}

极简操作、低学习成本：主界面仅保留「拍照识别」「搜索」「预警」三大入口；模型端侧离线运行，弱网、无网环境同样可用。

\subsection{落地性强}

贴合小农户生产习惯：术语白话化、方案本地化、用药合规化，避免「正确的废话」，输出可直接执行的行动建议。

\subsection{精准时效}

本地化、窗口期导向：方案匹配结合当地用药登记目录与季节窗口，预警推送基于 LBS 与气象、发病数据，保证「在对的时间、给对的建议」。

\section{技术架构与核心能力}

\subsection{轻量化 AI 图像识别模型}

采用深度可分离卷积神经网络，配合知识蒸馏与 INT8 量化，将模型体积压缩至数 MB，实现中低端手机离线秒级识别（详见第七章技术研发计划）。

\subsection{权威病虫害防治知识库}

联合农业院校与植保专家，构建覆盖主粮与经济作物的病虫害防治知识库，包含「病—药—量—时」结构化字段，并持续更新。

\subsection{低门槛适配能力}

支持主流安卓机型与中低端配置，对拍照质量、光线条件具备较强鲁棒性；对偏远地区网络条件不敏感。

\section{产品迭代规划}

\begin{table}[H]
  \centering
  \small
  \renewcommand{\arraystretch}{1.25}
  \caption{产品迭代路线图}
  \begin{tabular}{cll}
    \toprule
    阶段 & 时间 & 核心交付 \\
    \midrule
    MVP & 第 1--6 个月 & 水稻/小麦 30 类病虫害识别 + 基础防治方案 + 端侧部署 \\
    V1.0 & 第 7--12 个月 & 扩展至 60 类、知识库搜索、LBS 预警、县域试点 \\
    V2.0 & 第 13--24 个月 & 100+ 类、田间复杂场景目标检测、社群互助、多作物覆盖 \\
    V3.0 & 第 25--36 个月 & 数据飞轮、区域病害地图、农资闭环、开放 API \\
    \bottomrule
  \end{tabular}
\end{table}
